
\begin{obeassessment}
\textbf{Kuis Bab 11}
1. Jelaskan langkah-langkah matematis yang terjadi dalam pertukaran kunci Diffie-Hellman.
2. Apa yang dimaksud dengan masalah logaritma diskrit dan mengapa hal ini menjadi fondasi keamanan DH?
3. Gambarkan alur serangan \textit{Man-in-the-Middle} pada protokol DH dan jelaskan mengapa hal tersebut bisa terjadi.
\end{obeassessment}
\begin{competencychecklist}
\begin{tabularx}{\linewidth}{|X|c|}
\hline
Indikator Kompetensi & Tercapai \\
\hline
Saya dapat menjelaskan masalah distribusi kunci pada simetri & [ ] \\
\hline
Saya dapat menguraikan prosedur pertukaran kunci Diffie-Hellman & [ ] \\
\hline
Saya dapat menghitung kunci rahasia bersama $K$ menggunakan parameter DH & [ ] \\
\hline
Saya dapat menganalisis risiko serangan MITM pada protokol DH & [ ] \\
\hline
Saya dapat menjelaskan implementasi DH pada protokol TLS/SSL & [ ] \\
\hline
\end{tabularx}
\end{competencychecklist}
